\section{The middle-range error at a good point}

% BEGIN C78B1
\paragraph{C78B1. The declared row and column costs.}
Fix a datum $D$, its size data $L$, a hub $X$ and a twisted
localization $i$ with $\ell=\ell_i$ outside the two host primes
$g$ and satisfying the middle-range guard
$\mathrm{RangeAD}(L,X.\mathrm{exc},\ell)$. In particular
$\ell$ lies outside the small primes and differs from $q,d$.
Write $R_i,C_i$ for its declared sets of row and column places,
$d_u=e_uf_u$ for the local degree, $k=|C_i|$, $V=n_V$,
$H=\mathrm{large}_2(L)$, $v=v_{\rm del}(i)$, and
$h=h_0(i,C_i)$. For either set $S=R_i,C_i$ one has
\[
       |S|\le\sum_{u\in S}d_u\le [E:\mathbb Q]=2f .
\]
Indeed each $d_u\ge1$, and the declaration places all of $S$ over
$\ell$. The sum of ramification indices times residue degrees over
all primes over $\ell$ is the field degree. Therefore
$A=\sum_{u\in R_i}(d_u+2)\le 6f$ and $k\le2f$.
Moreover $0\le v\le 2t$: its sum over $0<a<t$ is bounded by
$t s_0(i)+v_\ell(t!)$, $s_0(i)=1$ off the small primes, and counting
multiples of $\ell^b$ gives $v_\ell(t!)\le t$.

Let $c\ge1$, and assume $b=X.c_B(i)\le ctnVH$ and $h\le ct$.
Put $M=cfn t^2VH$. The declaration part of the error is
\[
      t A b + 4V(k+1)v + t k h\le (6+24+2)M=32M .
\]
For the first term substitute $A\le6f$ and the bound for $b$.
For the second use $k+1\le3f$ and $v\le2t$. The third is bounded
by $2cft^2$. In each comparison $f,n,t,V,H\ge1$.
At a split prime and good $\pi$, suppose in addition that, for every
$u\in C_i$, $c_0(u)=X.c_0(i,u)\le cnVH$ and
$h_W(u)=h^0_W(i,u,X.\rho_i,\pi)\le crnVH$. The additional error is
\[
   \sum_{u\in C_i}\bigl(d_u m v +2v+h_W(u)+r c_0(u)\bigr)\le16 M .
\]
The first sum is at most $2f\cdot t\cdot2t\le4M$, since $m\le t$.
The second is at most $2f\cdot4t\le8M$. The last two together are
bounded by $2f\cdot2crnVH\le4M$, since $r\le t$. Empty sets
give the same inequalities by the empty-sum convention.
% END C78B1

% BEGIN C78B2
\paragraph{C78B2. From the coarse invariant bounds to the error bound $7200$.}
Suppose $\mathrm{DataAD}(L,HW,X,R)$ holds and the three coarse
bounds of C78B1 hold with the same $c\ge1$ for every index and
point under their indicated guards. Here the final split bound for
$h_W$ concerns every good $\pi$; $c_0$ itself is independent of
that point. Then $\mathrm{BBLe}(L,X,60c)$ holds. In fact $\ell$
has a split or inert class after the three exclusions in C78B1.
The split equality of $\mathrm{DataAD}$ identifies the stored hub
error with the sum of the declaration and additional error, at most
$48M$. Its inert equality identifies the stored error with the
declaration alone, at most $32M$. Both are bounded by $60MH$ since
$H\ge1$ and $M\ge0$. This proves the precise inequality
\[
 X.\mathrm{errorBoundHub}(i,\pi)\le 60c\,fn t^2 n_V H^2 .
\]
In particular the native predicate
$\mathrm{InvariantsBoundsWithCoarseColumnConstant}(L,X,120)$ gives
$\mathrm{BBLe}(L,X,7200)$. The canonical hub of F78D2 has the
required $\mathrm{DataAD}$ by construction. Any common $C_s$ in F78D3
is at least $22365$, hence the bound at $7200$ implies its field
$\mathrm{BBLe}(L,X,C_s)$ on this same hub.
% END C78B2

% BEGIN C78B3
\paragraph{C78B3. Two of the coarse constants on the canonical hub.}
On the F78D2 hub, $H=L.\mathrm{sc}.L\ge1$,
$L.\mathrm{sc}.b=e^H$ and $N=4nn_V\ge1$. Put
$a=\lfloor2e^H\rfloor$, $K=\lfloor\log_\ell(2a)\rfloor$,
$K'=\lfloor\log_\ell(2(a+1))\rfloor$.
For a positive integer $z\le 2^b e^H$ and $\ell\ge2$, its floor
logarithm is at most $b+2H$: take natural logarithms of
$2^{\lfloor\log_\ell z\rfloor}\le z$ and use $\log 2>1/2$.
Thus $K\le2+2H\le4H$ and $K'\le3+2H\le5H$.
The canonical definitions at these exact indices are
\[
 \begin{split}
 c_B&=6+K+7t+4tN+2tNK,\\
 c_0(u)&=5+K'+N(4+2K').
 \end{split}
\]
Since $t,N,H\ge1$, the five terms in $c_B$ are bounded by
$(6+4+7+4+8)tNH\le30tNH$; similarly
$c_0(u)\le(5+5+4+10)NH\le25NH$.
Substitute $N=4nn_V$. This gives $c_B\le120tnn_VH$ and
$c_0(u)\le100nn_VH\le120nn_VH$. Both hold on the original hub
at their full required ranges. No comparison of Galois
representations is needed to choose these constants.
% END C78B3

% BEGIN C78B4
\paragraph{C78B4. The scalar invariant term at every declared column.}
Under the middle-range guard the scalar term satisfies
$h_0(i,C_i)\le2t$. Write $O$ for the coefficient DVR and $K$ for
its fraction field. For $u\mid\ell$, the cyclotomic character takes
the value $1+\ell$ on a member of the local inertia group. To see
the relevant level, the defining host inclusion puts the local field
inside the compositum of the maximal unramified extension and
the $\ell^{s_0(i)}$-th roots of unity. Here $s_0(i)=1$ and $\ell$
is odd. On the cyclotomic tower take the automorphism of exponent
$1+\ell$ and the identity on the maximal unramified extension; the
cyclotomic extensions of $\mathbb Q_\ell$ are totally ramified,
as follows from the Eisenstein polynomial for a primitive root
minus $1$. Extend this compositum automorphism to the algebraic
closure. It fixes the host completion since it fixes $\mu_\ell$,
and its transport to the prescribed decomposition group is in inertia.
The bridge identifies its cyclotomic action on $O$ with $1+\ell$.

For $0<a<t$ an invariant in the source's inverse cyclotomic twist
of $K/O$ is therefore killed by $d=(1+\ell)^a-1$. For an odd prime
the binomial expansion shows that a power prime to $\ell$ leaves
$v_\ell(y-1)$ unchanged when $y\equiv1\pmod\ell$, and taking an
$\ell$-th power increases it by one. Thus
$v_O(d)=(1+v_\ell(a))v_O(\ell)$. The $d$-torsion of $K/O$ is
$d^{-1}O/O\cong O/(d)$, so its normalized length, and hence that
of the invariants, is at most $1+v_\ell(a)\le2t$. The last
inequality also follows from $2^{v_\ell(a)}\le a<t$.
The source takes the maximum over $(u,a)\in C_i\times(0,t)$
with value zero on the empty set. This proves the claimed bound,
in particular $h_0(i,C_i)\le120t$.
% END C78B4

% BEGIN C78B5
\paragraph{C78B5. The remaining annihilator on the split column.}
On the F78D2 hub, every guarded inert prime already satisfies the
$7200$ bound, by C78B1, C78B3 and C78B4 and the inert error identity
in C78B2. At split primes whose declared column is empty the same
argument applies, as the additional error sum vanishes.
For the remaining split primes there is a narrower sufficient
hypothesis for the entire $7200$ bound. At each guarded split $(i,u,\pi)$ in C78B1,
let $T$ be the standard rank-$r$ lattice of the localized twisted
representation defined from the source matrix $\mathrm{WpO}(i,X.\rho_i,\pi)$.
Write $V_T$ for its fraction space. It suffices to know, on the
prescribed local coefficient bridge at $u$, that
\[
     \ell^{c_0(u)} (V_T/T)^{G_{E_u}}=0 .
\]
In fact the coefficient bridge sends a vector class over the source DVR
to the coordinatewise class over the local DVR. This map is injective:
if all image coordinates are integral, pull them back through the
specified isomorphism of integer rings; injectivity on fraction fields
recovers the original coordinates. It intertwines the two matrix
actions by entrywise scalar extension. Thus a decomposition-invariant
source class has invariant image and is also annihilated by the stated
power. The $\ell^{d}$-torsion of $(K/O)^r$ for $d=c_0(u)$ is
$(\ell^{-d}O/O)^r$, of length $rd\,\mathrm{length}(O/\ell)$.
Dividing by $\mathrm{length}(O/\ell)$ gives
$h_W(u)\le r c_0(u)\le120rn n_VH$. Together with C78B3--C78B4
this supplies every coarse inequality in C78B2 on the same hub.

In the pinned route the annihilator uses the crystalline local lattice
at this $\pi$, its interval of jumps, the exclusion of a jump at zero
for every embedding, and the comparison of its calculated exponent
with the very $c_0(u)$ above. Once the invariant classes are killed
by $\ell^a$ for $a\le c_0(u)$, multiplication by
$\ell^{c_0(u)-a}$ proves the displayed hypothesis. This last
representation-theoretic annihilator is retained as an input here.
The arithmetic bounds for the canonical constants alone prescribe the
size of the exponent and do not establish its action on those classes.
% END C78B5
